% !TeX program = lualatex
\documentclass[b5paper, 10pt, abstract=true, captions=tableheading]{scrartcl}
\usepackage[english]{babel}

\usepackage[
  textwidth=136mm,
  textheight=205mm,
  headsep=10mm,
  top=18mm
]{geometry}
 
\usepackage{amsmath, mathtools}
\usepackage{amsthm, thmtools}
\usepackage{mdframed}
\usepackage{fontsetup}
\usepackage{setspace}
\usepackage{xcolor}

\setstretch{1.12}
\usepackage{microtype}
\usepackage{hyperref}
\usepackage{xcolor}
\hypersetup{linkbordercolor=red,
            urlbordercolor=blue,
            citebordercolor=blue,
            linkcolor=teal,
            urlcolor=cyan,
            citecolor=purple,
            colorlinks=true
}
\usepackage{cleveref}
\usepackage{enumerate}
\usepackage{graphicx}
\usepackage{adjustbox}
\usepackage{color}
\usepackage{booktabs}
\usepackage{tikz}
\usepackage{multirow}
\usepackage{cellspace}
\setlength{\cellspacetoplimit}{3pt}
\setlength{\cellspacebottomlimit}{3pt}
\usetikzlibrary{arrows.meta, fit, backgrounds, calc, decorations.pathreplacing}
\renewcommand\UrlFont{\color{blue}\rmfamily}
\urlstyle{rm}

\usepackage{csquotes}
\usepackage[backend=biber, sortcites=true, style=numeric-comp]{biblatex}
\addbibresource{refs.bib}
\theoremstyle{plain}
\newtheorem{theorem}{Theorem}
\newtheorem{lemma}{Lemma}
\newtheorem{corollary}{Corollary}
\newtheorem{proposition}{Proposition}
\newtheorem{remark}{Remark}
\newtheorem{claim}{Claim}
\theoremstyle{definition}
\newtheorem{definition}{Definition}
\declaretheorem[style=definition,qed=$\triangle$]{example}
\newtheorem{conjecture}{Conjecture}
\newtheorem{assumption}{Assumption}
\newtheorem{heuristic}{Heuristic}
\crefname{heuristic}{Heuristic}{Heuristics}
\crefname{assumption}{Assumption}{Assumptions}
\crefname{example}{Example}{Examples}


\setkomafont{section}{\normalfont\large\bfseries}
\renewcommand{\sectionformat}{\thesection\autodot\enskip}
\RedeclareSectionCommand[tocbeforeskip=0pt,tocentryformat=\normalfont\itshape,tocpagenumberformat=\normalfont,toclinefill=\TOCLineLeaderFill]{section}
\setkomafont{subsection}{\normalfont\normalsize\bfseries}
\renewcommand{\subsectionformat}{\thesubsection\autodot\enskip}
\setkomafont{paragraph}{\normalfont\bfseries}
\RedeclareSectionCommand[afterskip=-.5em]{paragraph}
%\setkomafont{title}{\large\bfseries}   % the title itself
\setkomafont{title}{\bfseries\boldmath}
\addtokomafont{title}{\let\huge\Large}
\setkomafont{subtitle}{\large}          % subtitle, if you use one
\setkomafont{author}{\normalsize}       % author name(s)
\setkomafont{date}{\small} 
\addtokomafont{disposition}{\rmfamily}


\usepackage{dsfont, stmaryrd}
\usepackage{xfrac, faktor}
\newcommand{\FF}{\mathds{F}}
\newcommand{\EE}{\mathds{E}}
\newcommand{\KK}{\mathds{K}}
\newcommand{\ZZ}{\mathds{Z}}
\newcommand{\NN}{\mathds{N}}
\newcommand{\GL}{\operatorname{GL}}
\newcommand{\tr}{\mathsf{T}}
\newcommand{\CC}{\mathds{C}}
\newcommand{\pub}{\texttt{pub}}
\newcommand{\supp}{\operatorname{supp}}
\newcommand{\rank}{\operatorname{rank}}
\newcommand{\PMat}{\operatorname{PMat}}
\newcommand{\Pct}{\operatorname{Pct}}
\newcommand{\Sh}{\operatorname{Sh}}
\newcommand{\RowSpan}{\operatorname{RowSpan}}
\newcommand{\boldalpha}{\symbf{\alpha}}
\newcommand{\boldlambda}{\symbf{\lambda}}
\newcommand{\boldbeta}{\symbf{\beta}}
\newcommand{\intrange}[2]{\llbracket #1,#2 \rrbracket}
\renewcommand{\vec}[1]{\symbf{#1}}
\newcommand{\mat}[1]{\symbf{#1}}
\newcommand{\GRS}{\operatorname{GRS}}
\newcommand{\RS}{\operatorname{RS}}
\newcommand{\diag}{\operatorname{diag}}
\newcommand{\calK}{\mathcal{K}}
\newcommand{\poly}{\operatorname{poly}}
\newcommand{\Span}{\operatorname{span}}
\newcommand{\Gal}{\operatorname{Gal}}
\newcommand{\bits}{\FF_2}
\newcommand{\Tr}{Tr}
\newcommand{\boldgamma}{\vec{\gamma}}
\newcommand{\Goppa}{\Gamma}
\newcommand{\lin}{\mathsf{lin}}
\newcommand{\id}{\operatorname{id}}
\newcommand{\sample}{\in_{\$}}
\newcommand{\corrterm}{\mathsf{corr}}
\newcommand{\orth}{orth}
\newcommand{\rem}{rem}
\newcommand{\streamgen}{\mathsf{G}}
\newcommand{\Hash}{\mathsf{H}}
\newcommand{\negl}{\operatorname{negl}}
\newcommand{\mxi}{{\mathsf{m}_{\xi}}}
\newcommand{\Adv}{\mathsf{Adv}}
\newcommand{\msg}{\mathsf{msg}}
\newcommand{\hasse}[1]{\partial^{[#1]}}
\newcommand{\jet}{J}       
\newcommand{\Idx}{\mathcal{I}} 
\newcommand{\Phom}{\mathcal{P}^{\mathsf{hom}}}
\newcommand{\Pred}{\mathcal{P}^{\mathsf{red}}}
\newcommand{\RM}{\operatorname{RM}}
\newcommand{\Cond}{\operatorname{Cond}}


\newcommand{\mult}{\operatorname{mult}}
\newcommand*\inst[1]{\textsuperscript{#1}}
\newcommand*\email[1]{\texttt{#1}}
\setkomafont{publishers}{\small} 


\begin{document}
\title{Sampling Binary Kernels with Block Wiedemann}
\subtitle{a companion note}
\date{\today}
\author{%
  Kirill Vedenev\thanks{%
    \begin{tabular}[t]{@{}l@{}}
      CryptoPro, Moscow, Russia.\\
      Southern Federal University, Rostov-on-Don, Russia.\\
      \email{vedenevk@yandex.ru}
    \end{tabular}%
  }%
}
\publishers{}
\maketitle


\begin{abstract}
We give a block Wiedemann algorithm for sampling independent uniform
elements of the kernel of a binary matrix. Random Toeplitz compression
and auxiliary columns yield an explicit rejection bound for every fixed
input matrix. Accepted outputs are exactly uniform, and the algorithm
uses the input matrix only through matrix-vector products. We apply the
sampler to the polynomial hold-out systems used in algebraic key recovery
for Classic McEliece. For five parameter sets, we give explicit Boolean
circuits, storage bounds, and schedules charging for memory access. One
attempt costs at most $2^{92.68}$ to $2^{101.21}$ Boolean gates, depending
on the parameter set, and rejects with probability below $2^{-62}$.
\end{abstract}

\setcounter{tocdepth}{1}
{\hypersetup{linkcolor=black}\tableofcontents}

\section{The kernel sampler}
\label{app:wiedemann}
Let $\mat K\in\FF_2^{J_{\calK}\times N}$ be a binary matrix and write
$\calK=\ker\mat K$. We seek $M$ independent uniform elements of $\calK$.
The sampled vectors may be zero or linearly dependent. We assume that
products $\mat K\vec x$ can be computed without storing the matrix.

We give a block Wiedemann sampler
\cite{Wiedemann1986,Coppersmith1994} with a proved rejection bound for
every fixed $\mat K$. Conditional on acceptance, its outputs have the
required uniform distribution. The construction and proof use only
matrix products and do not require geometric hypotheses on the kernel.

Our application is the polynomial hold-out kernel used for algebraic key
recovery in \cite{Vedenev2026KeyRecovery}. Its entries are derivatives of
monomials evaluated at public columns, so matrix products can be computed
from those columns. \Cref{app:wiedemann-cost} defines this matrix and
counts the sampler's gates, storage, and memory access at the Classic
McEliece parameters. These are sampling costs; the algebraic conditions
and costs of subsequent key recovery are studied in
\cite{Vedenev2026KeyRecovery}.

\subsection{The scalar algorithm}
\label{subsec:wiedemann-scalar}
Wiedemann's algorithm solves a square linear system using only products of
the matrix with vectors \cite{Wiedemann1986}. It obtains a linear recurrence
from the Berlekamp--Massey algorithm, which is familiar from the decoding of
BCH and Reed--Solomon codes \cite{Massey1969}.

Let $\mat A$ be a square matrix of order $N'$ over a field $\FF$, and let
$\vec b\in\FF^{N'}$. In our application, $\mat A$ will be built from
$\mat K$ in \Cref{subsec:wiedemann-sampler}, with $N'$ slightly larger than
$N$. The vectors
$\vec b,\mat A\vec b,\mat A^2\vec b,\ldots$ satisfy linear relations, since
at most $N'$ of them are linearly independent. The \emph{minimal
polynomial} of $\vec b$ with respect to $\mat A$ is the monic polynomial
\[
  \Theta(\xi)=\sum_{e=0}^{\deg\Theta}\theta_e\xi^e
\]
of least degree with $\Theta(\mat A)\vec b=\vec0$. Its degree is at most
$N'$. If $\theta_0\neq0$, this relation can be solved for $\vec b$:
\[
  \vec b=\mat A\vec z,
  \qquad
  \vec z=-\theta_0^{-1}\sum_{e=1}^{\deg\Theta}\theta_e\mat A^{e-1}\vec b.
\]
Once the coefficients of $\Theta$ are known, Horner's rule computes $\vec z$
with $\deg\Theta-1$ products with $\mat A$.

A solver gives kernel vectors. If $\mat A\vec z=\mat A\vec u$, then
$\vec u-\vec z\in\ker\mat A$. If moreover $\vec u$ is uniformly random and
$\vec z$ is computed from $\mat A\vec u$ alone, then $\vec u-\vec z$ is a
uniform element of $\ker\mat A$: given $\mat A\vec u$, the vector $\vec u$ is uniform on a
coset of $\ker\mat A$, and subtracting the fixed $\vec z$ maps this coset
onto $\ker\mat A$ (\Cref{lem:wiedemann-uniform}).

The formula for $\vec z$ needs $\theta_0\neq0$, that is,
$\vec b\in\Span\{\mat A\vec b,\mat A^2\vec b,\ldots\}$. This can fail for
the right-hand sides $\vec b=\mat A\vec u$ of the kernel sampler. For
instance, if $\vec b=\mat A\vec u\neq\vec0$ and $\mat A\vec b=\vec0$, then
$\Theta(\xi)=\xi$. The sampler of \Cref{subsec:wiedemann-sampler} adds
helper columns to handle this case.

The algorithm finds $\Theta$ from scalars, because the vectors
$\mat A^e\vec b$ are too many to store. It chooses a random vector
$\vec x\in\FF^{N'}$ and computes the sequence
\[
  a_e=\vec x^\tr\mat A^e\vec b
  \quad\text{for }0\leq e<2N'.
\]
The coefficients of $\Theta$ give a linear recurrence for this sequence:
\[
  \sum_{e=0}^{\deg\Theta}\theta_e\,a_{e+e'}
  =\vec x^\tr\mat A^{e'}\Theta(\mat A)\vec b=0
  \quad\text{for all }e'\geq0.
\]
The minimal polynomial of the sequence is the monic polynomial of least
degree whose coefficients give such a recurrence. Since it has degree at
most $N'$, the Berlekamp--Massey algorithm computes it from the first $2N'$
terms with $O(N'^2)$ operations in $\FF$. (In decoding, the algorithm
returns the connection polynomial, which is the reciprocal of this
polynomial.) The minimal polynomial of the sequence divides $\Theta$.

The two polynomials can differ, and over $\FF_2$ this can happen with
large probability. For instance, let $\mat A^2\vec b=\vec b\neq\mat A\vec b$, so that
$\Theta(\xi)=\xi^2+1$. The sequence alternates between $\vec x^\tr\vec b$
and $\vec x^\tr\mat A\vec b$. Its minimal polynomial is $\Theta$ if and only
if these two values differ, which has probability $1/2$ for uniform
$\vec x\in\FF_2^{N'}$. Otherwise the algorithm finds $\xi+1$ or $1$. With
$\xi+1$, the formula gives $\vec z=\vec b$, which is not a solution.

One run of the scalar algorithm uses about $2N'$ products to compute the
sequence and fewer than $N'$ products to compute $\vec z$. It stores only a
few vectors of length $N'$ at a time. Two properties make it unsuitable at
our sizes. Over $\FF_2$, the recurrence found from one random $\vec x$ can
be a proper divisor of $\Theta$, as above. The $2N'$ products also form one
chain, in which each product needs the previous one, so they cannot run in
parallel. The block algorithm of \Cref{subsec:wiedemann-block} replaces
$\vec x$ by several random rows, which makes it unlikely that a factor is
lost. It replaces $\vec b$ by several columns, which split the products into
independent chains.

\subsection{The block algorithm}
\label{subsec:wiedemann-block}
The block algorithm of Coppersmith \cite{Coppersmith1994} computes the
scalar sequences for many pairs of vectors at once and finds one relation
for all of them. From now on the field is $\FF_2$. The algorithm uses a
matrix $\mat V\in\FF_2^{N'\times b}$ with $b$ columns and a matrix
$\mat X\in\FF_2^{b'\times N'}$ with $b'$ rows, and computes the matrix
sequence
\[
  \mat S_e=\mat X\mat A^e\mat V\in\FF_2^{b'\times b}
  \quad\text{for }e\geq0.
\]
At each position, the entries of $\mat S_0,\mat S_1,\ldots$ form a scalar
sequence as in \Cref{subsec:wiedemann-scalar}, for one row of $\mat X$ and
one column of $\mat V$.

The minimal polynomial is replaced by a relation with matrix
coefficients. Suppose that matrices $\mat C_0,\ldots,\mat C_L$ with $b$
rows satisfy
\begin{equation}\label{eq:wiedemann-vector-relation}
  \sum_{e=0}^{L}\mat A^e\mat V\mat C_e=\mat 0.
\end{equation}
Multiplying by $\mat X\mat A^{e'}$ gives a recurrence for the sequence,
as in the scalar case:
\begin{equation}\label{eq:wiedemann-sequence-relation}
  \sum_{e=0}^{L}\mat S_{e+e'}\mat C_e=\mat 0
  \quad\text{for all }e'\geq0.
\end{equation}
As in \Cref{subsec:wiedemann-scalar}, a relation whose coefficient
$\mat C_0$ has a prescribed form gives solutions of linear systems
(\Cref{subsec:wiedemann-sampler}).

The columns of the matrices $\mat A^e\mat V$, $e\geq0$, span the
\emph{Krylov space} of $\mat V$. The degree $L$ of the relation depends on
the $b$ columns of $\mat V$. The columns of $\mat A^e\mat V$, $0\leq e\leq L$,
are $b(L+1)$ vectors of length $N'$, so a nonzero relation exists as soon
as $b(L+1)>N'$. This count suggests a degree about $N'/b$, but does not
ensure the prescribed constant coefficient needed by the sampler.
\Cref{lem:wiedemann-preimages,lem:wiedemann-right} establish that stronger
claim for its random starting columns. The
number $L'$ of equations depends on the $b'$ rows of $\mat X$. A solution of
\eqref{eq:wiedemann-sequence-relation} for $0\leq e'<L'$ is, with high
probability, a solution of \eqref{eq:wiedemann-vector-relation} once $L'$
is about $N'/b'$ and $b'$ exceeds $b$ by a margin
(\Cref{lem:wiedemann-left}). The margin is needed because the Krylov space
can contain up to $2^{(b+1)\rho}$ vectors whose minimal polynomial has
degree $\rho$, and $\mat X$ must detect all of them. These equations involve only the terms
$\mat S_e$ with $e<L+L'$. Computing them takes about $N'(1+b/b')$ products
with $\mat A$, which form $b$ independent chains of length about $L+L'$.
\Cref{subsec:wiedemann-relation} computes the coefficients $\mat C_e$
from these terms.

\subsection{The sampler}
\label{subsec:wiedemann-sampler}
The sampler applies the block algorithm to the right-hand sides
$\mat A\mat U$ for a random matrix $\mat U$. It returns the columns of
$\mat U+\mat Z$, where $\mat A\mat Z=\mat A\mat U$; over $\FF_2$ this is
the vector $\vec u-\vec z$ of \Cref{subsec:wiedemann-scalar}. The sampler
adds two ingredients to the block algorithm. A random compression turns
$\mat K$ into a square matrix $\mat A$, which has the same kernel as the
padded matrix $\mat K_+$ below, except with small probability. Random helper columns in $\mat V$ supply the
solutions that the right-hand sides alone may miss.

The Wiedemann algorithm needs a square matrix, whereas $\mat K$ may be
rectangular. Over $\FF_2$, the square matrix $\mat K^\tr\mat K$ can
have a larger kernel than $\mat K$, because the image of $\mat K$ may
contain nonzero vectors orthogonal to all of it. We therefore compress
$\mat K$ by a random matrix. We first append $\zeta\geq0$ zero columns to
$\mat K$ and obtain
\[
  \mat K_+=[\mat K\mid\mat 0]\in\FF_2^{J_{\calK}\times N'},
  \qquad N'=N+\zeta,
  \qquad \ker\mat K_+=\calK\times\FF_2^{\zeta}.
\]
Let $\mat Q\in\FF_2^{N'\times J_{\calK}}$ be a Toeplitz matrix whose
$N'+J_{\calK}-1$ diagonals are independent uniform bits, and put
$\mat A=\mat Q\mat K_+$. Then $\ker\mat K_+\subseteq\ker\mat A$, with
equality unless $\mat Q$ maps a nonzero vector of the image of $\mat K_+$
to zero. Since $\mat A$ has $N'$ rows, at least $\zeta$ more than the
dimension of the image of $\mat K_+$, equality fails with probability less
than $2^{-\zeta}$ (\Cref{lem:wiedemann-compression}). A uniform matrix of
this size could not be stored, whereas $\mat Q$ is stored by its
diagonals. A product with $\mat A$ is a product
with $\mat K$ followed by a product with $\mat Q$, which is a polynomial
multiplication. The proofs use only one property of $\mat Q$: for every
fixed nonzero vector $\vec v$, the vector $\mat Q\vec v$ is uniform.
This property also keeps small, on average, the
dimension of $\operatorname{im}\mat A\cap\ker\mat A$, which the helper
columns must cover (\Cref{lem:wiedemann-compression}).

Let $M<b$, $b+2\leq b'$, and $L,L'\geq1$. One attempt of the sampler
proceeds as follows.
\begin{enumerate}[(1)]
  \item Draw $\mat Q$ and matrices $\mat U\in\FF_2^{N'\times M}$,
  $\mat W\in\FF_2^{N'\times(b-M)}$ and $\mat X\in\FF_2^{b'\times N'}$,
  independently and uniformly. Put $\mat A=\mat Q\mat K_+$ and
  $\mat V=[\mat A\mat U\mid\mat W]$. The columns of $\mat W$ are the
  helper columns.
  \item Compute $\mat S_e=\mat X\mat A^e\mat V$ for $0\leq e<L+L'$.
  \item Find $\mat C_1,\ldots,\mat C_L\in\FF_2^{b\times M}$ that satisfy
  \eqref{eq:wiedemann-sequence-relation} for $0\leq e'<L'$, with
  \[
    \mat C_0=\mat E=\begin{bmatrix}\mat I_M\\\mat 0\end{bmatrix},
  \]
  by the method of \Cref{subsec:wiedemann-relation}. If no such matrices
  exist, reject.
  \item Compute
  $\mat Z=\sum_{e=1}^{L}\mat A^{e-1}\mat V\mat C_e$ by Horner's rule.
  \item If $\mat K_+(\mat U+\mat Z)=\mat 0$, output the first $N$
  coordinates of the $M$ columns of $\mat U+\mat Z$. Otherwise, reject.
\end{enumerate}
The prescribed coefficient $\mat C_0=\mat E$ selects the right-hand sides,
since $\mat V\mat E=\mat A\mat U$. It plays the role of the condition
$\theta_0\neq0$ of \Cref{subsec:wiedemann-scalar}. The helper columns have
zero coefficient in $\mat C_0$, so they enter only the solution $\mat Z$.
If the coefficients satisfy the vector relation
\eqref{eq:wiedemann-vector-relation}, then $\mat A\mat U+\mat A\mat Z=\mat
0$, so the columns of $\mat U+\mat Z$ lie in $\ker\mat A$. Step (5) checks
against $\mat K_+$ rather than $\mat A$. Hence a wrong relation, or a
compression that loses the kernel, can only cause a rejection.

Steps (1) and (2) take $M+b(L+L'-1)$ products with $\mat A$, step (4)
takes $M(L-1)$, and step (5) takes $M$ products with $\mat K_+$.

\begin{theorem}\label{thm:wiedemann-sampler}
Let $\mat K\in\FF_2^{J_{\calK}\times N}$ be any matrix, and let
$\zeta,M,b,b',L,L'$ be as above. One attempt of the sampler rejects with
probability less than
\[
  \varepsilon
  =2^{-\zeta}+2^{M-b}+2^{N'-bL}+2^{N'-b'L'}
   +\frac1{2^{b'-b-1}-1}.
\]
Whenever the attempt accepts with positive probability, its output
conditional on acceptance consists of $M$ independent uniform elements of
$\calK$.
\end{theorem}

Each term of $\varepsilon$ bounds the failure of one event of the proof.
The term $2^{-\zeta}$ bounds the loss of the kernel in the compression, and
$2^{M-b}$ bounds the failure of the $b-M$ helper columns. The terms
$2^{N'-bL}$ and $2^{N'-b'L'}$ decrease with the excess of $bL$ and $b'L'$
over $N'$. The last term bounds the failure of $\mat X$ on the vectors of
the Krylov space whose minimal polynomial has small degree, and it requires
$b'\geq b+2$ (\Cref{lem:wiedemann-left}). For an integer $\eta\geq0$, the choices
$L=\lceil(N'+\eta)/b\rceil$ and $L'=\lceil(N'+\eta)/b'\rceil$ make
$2^{N'-bL}$ and $2^{N'-b'L'}$ at most $2^{-\eta}$. Attempts with fresh
randomness are independent. When $\varepsilon<1$, the expected number of
attempts is less than $1/(1-\varepsilon)$, and every accepted output has
the distribution of the theorem. \Cref{app:wiedemann-cost} chooses the
parameters.
\subsection{Finding the relation}
\label{subsec:wiedemann-relation}
Step (3) is a linear system over $\FF_2$. For each column of $\mat E$, it
has $b'L'$ equations in the $bL$ entries of the corresponding columns of
$\mat C_1,\ldots,\mat C_L$, and its matrix is block Hankel in the terms
$\mat S_e$. Gaussian elimination would take about $N'^3$ operations. The
system is also a Pad\'e-type approximation problem, of the same form as the
key equation of Reed--Solomon decoding. With fast polynomial
multiplication, approximant bases solve it exactly at a cost quasi-linear
in $L+L'$ \cite{GiorgiJeannerodVillard2003}.

Put $\mat S(\xi)=\sum_{e=0}^{L+L'-1}\mat S_e\xi^e$ and, for coefficients
$\mat C_0,\ldots,\mat C_L$, put
$\mat C^\circ(\xi)=\sum_{e=0}^{L}\mat C_e\xi^{L-e}$. The matrix
$\mat C^\circ$ is the reversal of the relation, as the connection
polynomial of \Cref{subsec:wiedemann-scalar} is the reciprocal of the
minimal polynomial. The coefficient of
$\xi^{L+e'}$ in $\mat S(\xi)\mat C^\circ(\xi)$ is
$\sum_{e}\mat S_{e+e'}\mat C_e$. Hence the equations of step (3) say that
the coefficients of $\xi^L,\ldots,\xi^{L+L'-1}$ in
$\mat S(\xi)\mat C^\circ(\xi)$ vanish, that is,
\begin{equation}\label{eq:wiedemann-pade}
  \mat S(\xi)\mat C^\circ(\xi)+\mat P(\xi)\equiv\mat 0
  \pmod{\xi^{L+L'}}
\end{equation}
for some polynomial matrix $\mat P$ of degree less than $L$. The
coefficient $\mat C_0$ is the coefficient of $\xi^L$ in $\mat C^\circ$.

We need all solutions of \eqref{eq:wiedemann-pade} within the two degree
bounds, so that we can impose the coefficient $\mat C_0=\mat E$ afterward.
An approximant basis gives a finite set of polynomial vectors from which
all these solutions can be formed.

Let $\mathcal A$ be the set of column vectors
$\vec v=(\vec c,\vec p)\in\FF_2[\xi]^{b+b'}$ with
$\mat S\vec c+\vec p\equiv\vec0\pmod{\xi^{L+L'}}$. It is a module over
$\FF_2[\xi]$ that contains $\xi^{L+L'}\FF_2[\xi]^{b+b'}$, so it is free of
rank $b+b'$. A basis of $\mathcal A$ is an \emph{approximant basis}. We
combine the two degree bounds into one by defining the \emph{shifted degree}
\[
  \operatorname{sdeg}\vec v=\max\{\deg\vec c,\ \deg\vec p+1\},
\]
where the degree of a polynomial vector is the largest degree of its
entries. Then \eqref{eq:wiedemann-pade} holds with $\deg\mat C^\circ\leq L$
and $\deg\mat P<L$ if and only if every column of
$(\mat C^\circ,\mat P)$ lies in $\mathcal A$ and has shifted degree at
most $L$.

The \emph{leading vector} $\operatorname{lv}(\vec v)\in
\FF_2^{b+b'}$ consists of the coefficients of $\xi^{\rho}$ in $\vec c$
and of $\xi^{\rho-1}$ in $\vec p$, where
$\rho=\operatorname{sdeg}\vec v$. A basis is \emph{reduced} if the
leading vectors of its elements are linearly independent. The algorithms
of \cite{BeckermannLabahn1994,GiorgiJeannerodVillard2003} compute a reduced
approximant basis from $\mat S_0,\ldots,\mat S_{L+L'-1}$, and
\cite{Thome2002} applies such algorithms to block Wiedemann.
\Cref{app:wiedemann-cost} implements the order-by-order construction
with fixed circuit addresses and counts its gates and storage.

Independence of the leading vectors prevents cancellation of the highest
shifted terms. It therefore bounds the coefficients in a combination of
basis vectors. The following lemma reduces the prescribed-coefficient
condition to a binary linear system.

\begin{lemma}\label{lem:wiedemann-predictable}
Let $\vec v_1,\ldots,\vec v_{b+b'}$ be a reduced approximant basis, with
$\rho_j=\operatorname{sdeg}\vec v_j$ and $\vec v_j=(\vec c_j,\vec p_j)$.
Let $\vec e_\nu$ be the column $\nu$ of $\mat E$.
Step (3) has a solution if and only if every column of $\mat E$ lies in
the span of the vectors $[\xi^{\rho_j}]\vec c_j$ with $\rho_j\leq L$.
If $\vec e_\nu=\sum_{\rho_j\leq L}\gamma_{\nu,j}[\xi^{\rho_j}]\vec c_j$ for
$1\leq\nu\leq M$, then the polynomial matrix $\mat C^\circ$ with columns
$\sum_{\rho_j\leq L}\gamma_{\nu,j}\xi^{L-\rho_j}\vec c_j$ gives a
solution.
\end{lemma}
\begin{proof}
Let $a_1,\ldots,a_{b+b'}$ be polynomials, not all zero, and put
$\rho=\max_j(\deg a_j+\rho_j)$. In $\vec v=\sum_ja_j\vec v_j$, the
coefficients of $\xi^\rho$ in $\vec c$ and of $\xi^{\rho-1}$ in
$\vec p$ form the vector $\sum_j\operatorname{lc}(a_j)\operatorname{lv}(\vec
v_j)$, summed over the $j$ with $\deg a_j+\rho_j=\rho$, where
$\operatorname{lc}$ is the leading coefficient. This vector is nonzero
because the basis is reduced, and no higher coefficients occur. Hence
$\operatorname{sdeg}\vec v=\rho$.

Every element of $\mathcal A$ is such a combination. By the first
paragraph, it has shifted degree at most $L$ if and only if
$\deg a_j\leq L-\rho_j$ for every $j$, which requires $a_j=0$ when
$\rho_j>L$. Its coefficient of $\xi^L$ in $\vec c$ is then
$\sum_{\rho_j\leq L}([\xi^{L-\rho_j}]a_j)\,[\xi^{\rho_j}]\vec c_j$.
These are therefore exactly the possible columns of $\mat C_0$, which
proves the first claim. For the second, the given columns are elements of
$\mathcal A$ of shifted degree at most $L$ whose coefficients of $\xi^L$
are the columns of $\mat E$.
\end{proof}

We use this method in step (3), with a fixed deterministic algorithm for
the basis and for the coefficients $\gamma_{\nu,j}$. It decides exactly
whether a solution exists, and the solution it returns is determined by
the terms $\mat S_e$. \Cref{lem:wiedemann-uniform} uses this.
\subsection{Proof of the sampler theorem}
\label{subsec:wiedemann-proofs}
Let $\mathcal V$ be the Krylov space of $\mat V$, and let $\mathcal V_L$ be
the span of the columns of $\mat A^e\mat V$ for $0\leq e<L$. The space
$\mathcal V$ is the smallest $\mat A$-invariant space that contains the
columns of $\mat V$. The proof shows that the attempt accepts when the
following four events occur, and bounds the probability that each of them
fails:
\begin{enumerate}[(i)]
  \item $\ker\mat A=\ker\mat K_+$;
  \item every column of $\mat A\mat U$ equals $\mat A\vec v$ for some
  $\vec v\in\mathcal V$;
  \item $\mathcal V_L=\mathcal V$;
  \item no nonzero $\vec v\in\mathcal V$ satisfies
  $\mat X\mat A^{e'}\vec v=\vec 0$ for all $0\leq e'<L'$.
\end{enumerate}
Uniformity is proved separately.

The first lemma concerns the compression and event (i). Let
$\chi=\dim(\operatorname{im}\mat A\cap\ker\mat A)$. This is the number of
Jordan blocks of size at least two for the eigenvalue $0$ of $\mat A$, and
it is the defect that the helper columns repair
(\Cref{lem:wiedemann-preimages}). The second claim of the lemma shows
that the random $\mat Q$ keeps $\chi$ small.

\begin{lemma}\label{lem:wiedemann-compression}
The event $\ker\mat A\neq\ker\mat K_+$ has probability less than
$2^{-\zeta}$. Moreover,
\[
  \mathbb E\bigl[(2^{\chi}-1)\cdot\mathbf 1\{\ker\mat A=\ker\mat K_+\}\bigr]
  <1.
\]
\end{lemma}
\begin{proof}
First we show that $\mat Q\vec v$ is uniform for every fixed nonzero
$\vec v\in\FF_2^{J_{\calK}}$. Write the entries of $\mat Q$ as
$Q_{l,l'}=\mu_{l-l'+J_{\calK}}$, where $\mu_1,\ldots,\mu_{N'+J_{\calK}-1}$ are
the diagonal bits, and let $l_0$ be the first index with $v_{l_0}=1$. The
coordinate $l$ of $\mat Q\vec v$ equals $\mu_{l-l_0+J_{\calK}}$ plus a sum
of bits $\mu_{l''}$ with $l''<l-l_0+J_{\calK}$. The bits
$\mu_{l-l_0+J_{\calK}}$, $1\leq l\leq N'$, are distinct. Hence, for fixed
values of the other bits, the map from these $N'$ bits to $\mat Q\vec v$
is a triangular bijection onto $\FF_2^{N'}$.

As noted in \Cref{subsec:wiedemann-sampler}, $\ker\mat A=\ker\mat K_+$
unless $\mat Q\vec v=\vec0$ for some nonzero $\vec v$ in the image of
$\mat K_+$. This image has fewer than $2^{N}$ nonzero vectors, and each is
mapped to zero with probability $2^{-N'}$. The union bound gives the first
claim.

For the second claim, let
$\mathcal Y=\{\vec v\in\operatorname{im}\mat K_+:\mat Q\vec v\in\ker\mat K_+\}$.
A fixed nonzero $\vec v\in\operatorname{im}\mat K_+$ lies in $\mathcal Y$
with probability $|\ker\mat K_+|/2^{N'}=2^{-\rank\mat K_+}$. Summing over
these vectors gives
\[
  \mathbb E\bigl[2^{\dim\mathcal Y}-1\bigr]
  =\bigl(2^{\rank\mat K_+}-1\bigr)2^{-\rank\mat K_+}<1.
\]
If $\ker\mat A=\ker\mat K_+$, then $\mat Q$ is injective on
$\operatorname{im}\mat K_+$, and
$\operatorname{im}\mat A\cap\ker\mat A=\mat Q(\mathcal Y)$. Hence
$\chi=\dim\mathcal Y$ on this event. Since $2^{\dim\mathcal Y}-1\geq0$,
the second claim follows.
\end{proof}

The second lemma concerns event (ii). It shows that the helper columns
supply solutions where the formula of \Cref{subsec:wiedemann-scalar}
fails.

\begin{lemma}\label{lem:wiedemann-preimages}
Fix $\mat A$ and $\mat U$. With probability at least
$1-(2^\chi-1)2^{M-b}$ over $\mat W$, every column of $\mat A\mat U$ equals
$\mat A\vec v$ for some $\vec v\in\mathcal V$.
\end{lemma}
\begin{proof}
On the part of $\FF_2^{N'}$ where $\mat A$ is invertible, every column of
$\mat A\mat U$ has a preimage in $\mathcal V$ without help. On the part
where $\mat A$ is nilpotent, the helper columns provide one. We first
split $\FF_2^{N'}$ into these two parts, by the Fitting decomposition.

The projections onto the two parts must preserve $\mathcal V$.
It is enough to express them as polynomials in $\mat A$, because
$\mathcal V$ is $\mat A$-invariant. B\'ezout's identity supplies them.
Write the minimal polynomial of $\mat A$ as $\xi^{e_0}g_1(\xi)$ with
$g_1(0)=1$. Since $\xi^{e_0}$ and $g_1$ are coprime, there are polynomials
$g_2,g_3$ with $g_2(\xi)\xi^{e_0}+g_3(\xi)g_1(\xi)=1$. The matrices
$\mat\Pi_0=g_3(\mat A)g_1(\mat A)$ and $\mat\Pi_1=g_2(\mat A)\mat A^{e_0}$
satisfy $\mat\Pi_0+\mat\Pi_1=\mat I$ and
$\mat A^{e_0}\mat\Pi_0=g_1(\mat A)\mat\Pi_1=\mat 0$. Moreover, $\mat\Pi_1$
vanishes on $\ker\mat A^{e_0}$ and $\mat\Pi_0$ vanishes on
$\ker g_1(\mat A)$. Hence they are the projections onto $\mathcal F_0=\ker\mat A^{e_0}$ and
$\mathcal F_1=\ker g_1(\mat A)$ in the decomposition
$\FF_2^{N'}=\mathcal F_0\oplus\mathcal F_1$. Both spaces are
$\mat A$-invariant, $\mat A$ is nilpotent on $\mathcal F_0$, and
$\ker\mat A\subseteq\mathcal F_0$. On $\mathcal F_1$, the identity
$g_1(\mat A)=\mat 0$ gives $\mat A^{-1}=g_4(\mat A)$ with
$g_4(\xi)=(g_1(\xi)+1)/\xi$.

Let $\vec u$ be a column of $\mat U$, and write
$\vec u=\vec u_0+\vec u_1$ with $\vec u_0\in\mathcal F_0$ and
$\vec u_1\in\mathcal F_1$. The space $\mathcal V$ contains $\mat A\vec u$
and is invariant under every polynomial in $\mat A$. On the invertible
part, this gives
$\vec u_1=g_4(\mat A)\mat\Pi_1\mat A\vec u\in\mathcal V$, for every choice
of $\mat W$.

On the nilpotent part, two vectors give the same image under $\mat A$
exactly when they differ by an element of $\ker\mat A$. We therefore work
in the quotient $\overline{\mathcal F}_0=\mathcal F_0/\ker\mat A$.
The induced action of $\mat A$ is still nilpotent. If $e_0=0$, then $\mathcal F_0=\{\vec 0\}$ and there is
nothing to prove, so let $e_0\geq1$. It suffices that the image
$\mathcal F'$ of $\mathcal V\cap\mathcal F_0$ in $\overline{\mathcal F}_0$
is all of $\overline{\mathcal F}_0$. Indeed, there is then a
$\vec v_0\in\mathcal V\cap\mathcal F_0$ with
$\vec u_0-\vec v_0\in\ker\mat A$, and $\vec v=\vec v_0+\vec u_1\in\mathcal
V$ satisfies $\mat A\vec v=\mat A\vec u$. The image $\mathcal F'$ is
$\mat A$-invariant, and $\mat A$ is nilpotent on
$\overline{\mathcal F}_0$. It therefore suffices that $\mathcal F'$ spans
$\overline{\mathcal F}_0$ modulo $\mat A\overline{\mathcal F}_0$. Indeed,
if $\mathcal F'+\mat A\overline{\mathcal F}_0=\overline{\mathcal F}_0$,
then substituting this identity into itself and using
$\mat A\mathcal F'\subseteq\mathcal F'$ gives
$\overline{\mathcal F}_0=\mathcal F'+\mat A^e\overline{\mathcal F}_0$ for
every $e\geq1$, and $\mat A^{e_0}\overline{\mathcal F}_0=\{\vec 0\}$.

The right-hand sides cannot do this. The component of $\mat A\vec u$ in
$\mathcal F_0$ is $\mat A\vec u_0\in\mat A\mathcal F_0$, so its class in
$\overline{\mathcal F}_0/\mat A\overline{\mathcal F}_0$ is zero. The
helper columns must span this quotient. We first show that
$\dim(\overline{\mathcal F}_0/\mat A\overline{\mathcal F}_0)=\chi$. This
quotient is $\mathcal F_0/(\ker\mat A+\mat A\mathcal F_0)$. Since
$\dim\ker\mat A+\dim\mat A\mathcal F_0=\dim\mathcal F_0$, its dimension is
$\dim(\ker\mat A\cap\mat A\mathcal F_0)$. This intersection equals
$\ker\mat A\cap\operatorname{im}\mat A$. Indeed, let
$\mat A\vec v'\in\ker\mat A$ with $\vec v'=\vec v'_0+\vec v'_1$,
$\vec v'_0\in\mathcal F_0$ and $\vec v'_1\in\mathcal F_1$. Then
$\mat A\vec v'$ and $\mat A\vec v'_0$ lie in $\mathcal F_0$. Hence
$\mat A\vec v'_1\in\mathcal F_1\cap\mathcal F_0=\{\vec 0\}$, so
$\mat A\vec v'=\mat A\vec v'_0\in\mat A\mathcal F_0$.

The projections by $\mat\Pi_0$ of the columns of $\mat W$ lie in
$\mathcal V\cap\mathcal F_0$. The map that sends a vector to the class of
its projection in
$\overline{\mathcal F}_0/\mat A\overline{\mathcal F}_0$ is linear and
surjective. Hence the $b-M$ classes are independent uniform vectors of a
space of dimension $\chi$. They fail to span it only if some nonzero
linear functional vanishes on all of them. Each such functional vanishes
with probability $2^{M-b}$, and there are $2^\chi-1$ choices. The union
bound gives failure probability at most $(2^\chi-1)2^{M-b}$.
\end{proof}

The third lemma concerns event (iii): the first $L$ matrices
$\mat A^e\mat V$ already span the Krylov space. This makes the degree
count of \Cref{subsec:wiedemann-block} precise. A row vector $\vec x^\tr$
has a minimal polynomial with respect to $\mat A$, the monic polynomial
$\Theta$ of least degree with $\vec x^\tr\Theta(\mat A)=\vec 0$.

\begin{lemma}\label{lem:wiedemann-right}
For every fixed $\mat A$, the event $\mathcal V_L\neq\mathcal V$ has
probability less than $2^{N'-bL}$ over $\mat U$ and $\mat W$.
\end{lemma}
\begin{proof}
The space $\mathcal V_L$ equals $\mathcal V$ if and only if it is
$\mat A$-invariant, that is, if and only if it contains the columns of
$\mat A^L\mat V$. If it does not, some nonzero row vector $\vec x^\tr$
vanishes on $\mathcal V_L$ but not on all columns of $\mat A^L\mat V$. The
minimal polynomial of such an $\vec x^\tr$ has degree greater than $L$.
Otherwise it would express $\vec x^\tr\mat A^L$ as a combination of
$\vec x^\tr\mat A^{e}$ with $e<L$, and $\vec x^\tr\mat A^L$ would vanish
on the columns of $\mat V$.

Fix a nonzero row vector $\vec x^\tr$ whose minimal polynomial has degree
greater than $L$. The rows
$\vec x^\tr,\vec x^\tr\mat A,\ldots,\vec x^\tr\mat A^L$ are then linearly
independent. The row $\vec x^\tr$ vanishes on $\mathcal V_L$ if and only if
every column of $\mat W$ is orthogonal to $\vec x^\tr\mat A^e$ for
$0\leq e<L$, and every column of $\mat U$ is orthogonal to
$\vec x^\tr\mat A^e$ for $1\leq e\leq L$. Each of these
$b$ columns must satisfy $L$ independent linear conditions, which has
probability $2^{-L}$ per column, and $2^{-bL}$ for all $b$ columns
together, since they are independent. The union bound over fewer than $2^{N'}$
row vectors gives the claim.
\end{proof}

The fourth lemma concerns event (iv): the projection by $\mat X$ loses no
information about $\mathcal V$.

\begin{lemma}\label{lem:wiedemann-left}
For every fixed $\mat A$ and $\mat V$, the probability over $\mat X$ that
some nonzero $\vec v\in\mathcal V$ satisfies
$\mat X\mat A^{e'}\vec v=\vec 0$ for $0\leq e'<L'$ is less than
\[
  2^{N'-b'L'}+\frac1{2^{b'-b-1}-1}.
\]
\end{lemma}
\begin{proof}
A vector with a low-degree minimal polynomial gives fewer independent
conditions on the random projection. We must count these vectors more
carefully than by the size of the ambient space.

First fix a polynomial that annihilates the vector.
The space $\mathcal V$ is a module over $\FF_2[\xi]$, with $\xi$ acting as
$\mat A$, generated by the $b$ columns of $\mat V$. Let $\Theta$ be a
polynomial of degree $\rho$. The quotient
$\mathcal V/\Theta(\mat A)\mathcal V$ is generated by $b$ elements over
$\FF_2[\xi]/(\Theta)$, so its dimension is at most $b\rho$. Since
$\Theta(\mat A)$ is an endomorphism of $\mathcal V$, its kernel has the
same dimension as this cokernel. Hence at most $2^{b\rho}$ vectors of
$\mathcal V$ satisfy $\Theta(\mat A)\vec v=\vec 0$.

Let $\vec v\in\mathcal V$ be nonzero, with minimal polynomial of degree
$\rho$, and put $\rho'=\min(\rho,L')$. The vectors $\mat A^e\vec v$ with
$0\leq e<\rho'$ are linearly independent. Each $\mat A^e\vec v$ with
$e\geq\rho$ is a combination of those with $e<\rho$, so the conditions
$\mat X\mat A^{e'}\vec v=\vec 0$ for $e'<L'$ are equivalent to the same
conditions for $e'<\rho'$. Since $\mat X$ is uniform, the products of
$\mat X$ with these $\rho'$ vectors are independent uniform vectors of
$\FF_2^{b'}$. All of them vanish with probability $2^{-b'\rho'}$.

For $\rho<L'$, there are $2^\rho$ monic polynomials of degree $\rho$, and
each is the minimal polynomial of at most $2^{b\rho}$ vectors. For
$\rho\geq L'$, we use the bound $2^{N'}$ on the number of vectors. The
union bound gives
\[
  2^{N'}2^{-b'L'}+\sum_{\rho=1}^{L'-1}2^{\rho+b\rho-b'\rho}
  <2^{N'-b'L'}+\frac1{2^{b'-b-1}-1}.
  \qedhere
\]
\end{proof}

The last lemma proves uniformity. It uses no probability bound of the
preceding lemmas.

\begin{lemma}\label{lem:wiedemann-uniform}
Conditional on acceptance, the columns of $\mat U+\mat Z$ are independent
uniform elements of $\ker\mat K_+$.
\end{lemma}
\begin{proof}
The attempt uses $\mat U$ only through $\mat K_+\mat U$. Indeed,
$\mat A\mat U=\mat Q(\mat K_+\mat U)$ determines $\mat V$ and the
sequence. By \Cref{subsec:wiedemann-relation}, the sequence determines
the coefficients $\mat C_e$. These determine $\mat Z$, and the check of
step (5) compares $\mat K_+\mat Z$ with $\mat K_+\mat U$. Fix the
value of $\mat K_+\mat U$ together with $\mat Q$, $\mat W$ and $\mat X$.
Then $\mat Z$ and the acceptance are fixed, and $\mat U$ is uniform on the
set of matrices with this value of $\mat K_+\mat U$. This set is a translate of
$(\ker\mat K_+)^M$. Suppose that the attempt accepts. Then
$\mat K_+(\mat U+\mat Z)=\mat 0$, so $\mat U\mapsto\mat U+\mat Z$ maps the
set bijectively onto $(\ker\mat K_+)^M$. The output is therefore uniform
on $(\ker\mat K_+)^M$ for every fixed value of the other data, and hence
also after averaging over them.
\end{proof}

\begin{proof}[Proof of \Cref{thm:wiedemann-sampler}]
Suppose that the events (i)--(iv) occur. By (ii) and (iii), each column
$\mat A\vec u$ of $\mat A\mat U$ equals $\mat A\vec v$ with
$\vec v=\sum_{e<L}\mat A^e\mat V\vec c_e$ for some vectors
$\vec c_e\in\FF_2^b$. Collecting these vectors for the $M$ columns into
matrices gives $\mat A\mat U=\sum_{e=1}^{L}\mat A^e\mat V\mat C_e$ for some
$\mat C_e\in\FF_2^{b\times M}$. Together with $\mat C_0=\mat E$, they
satisfy \eqref{eq:wiedemann-vector-relation}, hence also
\eqref{eq:wiedemann-sequence-relation}. Step (3) therefore finds a
solution (\Cref{lem:wiedemann-predictable}) and does not reject.

Let $\mat C_1,\ldots,\mat C_L$ be the matrices found in step (3), and put
$\mat\Delta=\sum_{e=0}^{L}\mat A^e\mat V\mat C_e$. Its columns lie in
$\mathcal V$, and \eqref{eq:wiedemann-sequence-relation} for $e'<L'$ says
that $\mat X\mat A^{e'}\mat\Delta=\mat 0$. By (iv), $\mat\Delta=\mat 0$.
Since $\mat\Delta=\mat A\mat U+\mat A\mat Z$, the columns of $\mat U+\mat Z$ lie in
$\ker\mat A$. By (i), they lie in $\ker\mat K_+$, and step (5) accepts.

It remains to bound the probability that an event fails. The matrices
$\mat Q$, $\mat U$, $\mat W$ and $\mat X$ are independent, which is what
\Cref{lem:wiedemann-preimages,lem:wiedemann-right,lem:wiedemann-left}
use. Event (i) fails with probability less than $2^{-\zeta}$
(\Cref{lem:wiedemann-compression}). For fixed $\mat Q$ and $\mat U$, (ii)
fails with probability at most $(2^\chi-1)2^{M-b}$
(\Cref{lem:wiedemann-preimages}). Hence, by
\Cref{lem:wiedemann-compression},
\[
  \Pr\bigl[\text{(i) holds and (ii) fails}\bigr]
  \leq2^{M-b}\,\mathbb E\bigl[(2^\chi-1)\cdot\mathbf 1\{\text{(i)}\}\bigr]
  <2^{M-b}.
\]
Events (iii) and (iv) fail with probability less than
$2^{N'-bL}$ and $2^{N'-b'L'}+1/(2^{b'-b-1}-1)$
(\Cref{lem:wiedemann-right,lem:wiedemann-left}). The union bound gives
$\varepsilon$. Uniformity follows from \Cref{lem:wiedemann-uniform}, since
removing the last $\zeta$ coordinates maps a uniform element of
$\calK\times\FF_2^\zeta$ to a uniform element of $\calK$.
\end{proof}


\section{Circuit and access costs of the optimized sampler}
\label{app:wiedemann-cost}
The sampler of \Cref{app:wiedemann} obtains uniform kernel samples with
a proved bound on rejection. At the Classic McEliece parameters, the
implementation below has a gate count close to Weis's published block
Lanczos estimate. It costs more in the model that charges for memory
access. \Cref{tab:wopt-optimized-comparison} gives both comparisons.

Lowering selected multiplicities allows further shortening under the
kernel-dimension prediction. Wider projection blocks shorten the Krylov
sequences, at a cost in storage. Grouping the
derivative equations lets many rows share a matrix multiplication.
The Toeplitz products use a binary polynomial FFT, with operations on
known zero inputs and unused outputs removed. These choices leave the
sampling proof unchanged: it applies to every fixed matrix of equations.

We count gates of a fixed Boolean circuit. Every AND and XOR counts,
even when its control bit happens to be zero. The addresses and the
number of steps depend only on the parameters; public point coordinates
and fresh random bits are inputs. The access column also charges for the
size of the memory region used by each gate and each transfer.

\begin{table}[htbp]
\centering\small\setlength{\tabcolsep}{3pt}
\caption{One sampler attempt and Weis's published Lanczos estimates.
Gate and access columns are base-two logarithms. Our bounds are rounded
upward to $0.01$, and storage upward to $0.001$ TiB. The two constructions
use related hold-out systems, not the same matrix. Weis's figures are
from \cite[Tables 4--5 and Appendix E]{Weis2026}.}
\label{tab:wopt-optimized-comparison}
\begin{tabular}{@{}lrrrrr@{}}
\toprule
 &\multicolumn{2}{c}{Gates}&\multicolumn{2}{c}{Access}&\\
Parameter set & Weis & Ours & Weis & Ours & Our TiB\\
\midrule
\texttt{mceliece348864}&92.30&92.68&107.30&109.21&104.168\\
\texttt{mceliece460896}&98.34&98.19&113.96&115.21&851.026\\
\texttt{mceliece6688128}&101.31&101.21&117.20&119.06&2,322.558\\
\texttt{mceliece6960119}&100.55&100.35&116.37&117.68&1,791.580\\
\texttt{mceliece8192128}&101.31&101.21&117.20&119.06&2,322.558\\
\bottomrule
\end{tabular}
\end{table}

The gate counts differ by less than $0.4$ in the base-two logarithm.
Our access charges are about $2.4$ to $3.8$ times larger. The sampler
rejects with probability below $2^{-62}$ per attempt for the choices
below. Weis's estimated run length and supply of enough samples rely on
\cite[Assumption 4]{Weis2026}.

\providecommand{\ceil}[1]{\left\lceil #1\right\rceil}
\providecommand{\Cfive}{\binom{k}{5}}
\providecommand{\Cthree}{\binom{k}{3}}
\providecommand{\rhoo}[1]{\rho\!\left(#1\right)}
\subsection{The mixed hold-out system}
\label{sec:wopt-parameters}
The number of unknown coefficients is $\binom{k}{7}$. Reducing the
shortened dimension $k$ therefore reduces the coefficient vectors used
throughout the computation.
We lower the multiplicity at selected retained points to keep the
predicted kernel nonzero after further shortening. Following the
parameter selection in \cite[Sec.~9]{Vedenev2026KeyRecovery}, we impose
the minimum total multiplicity used by the vanishing criterion with the
$t$ hidden Goppa conditions.

Consider a shortened binary Goppa code over $\FF_2$, defined over
$\FF_{2^m}$ by a square-free Goppa polynomial of degree $t$. Its public
generator matrix is in systematic form, with columns
$\mathbf y_j\in\FF_2^k$. The length is
\(n_{\mathrm{sh}}=mt+k\), with polynomial degree bound
\(D_{\mathrm{sh}}=k+(m-2)t-1\). Choose one non-unit column as the held-out
position and omit it from the constraints. Include the public point at
infinity
\(\mathbf w=\sum_j\mathbf y_j\), at multiplicity six, in every parameter
set. There are then \(P=mt\) retained non-unit constraint points, counting
\(\mathbf w\). Of the finite retained points, choose \(P_5\) at
multiplicity five; the remaining \(P_6=P-P_5\) points have multiplicity six.
The \(k\) unit columns already have multiplicity at least six for every
homogeneous multilinear polynomial of degree seven.

The unknowns are the coefficients $h_I$ of
\[
 h(X_0,\ldots,X_{k-1})=\sum_{|I|=7}h_I\prod_{u\in I}X_u.
\]
For a set $S$ of coordinate indices, write $D_S^{[|S|]}h(y)$ for the
coefficient of $\prod_{u\in S}Z_u$ in $h(y+Z)$.
Let \(K\in\FF_2^{J\times N}\) map $(h_I)_{|I|=7}$ to these derivative
values at the retained points. At a point of multiplicity six, its rows
have $|S|=5$ and $3$; at a point of multiplicity five, they have $|S|=4$
and $3$. The entry in row $(y,S)$ and column $I$ is
$\prod_{u\in I\setminus S}y_u$ if $S\subseteq I$, and zero otherwise.
These are the mixed hold-out constraints of
\cite[Sec.~9.2]{Vedenev2026KeyRecovery}. Thus
\begin{equation}
 N=\binom{k}{7},\qquad
 J=P_6\binom{k}{5}+P_5\binom{k}{4}+P\binom{k}{3}.
 \label{eq:wopt-dimensions}
\end{equation}
\Cref{tab:wopt-parameters} lists the choices used in the cost analysis.
We take $M=k-m+31$, the sample count derived for this application in
\cite[Sec.~9.1]{Vedenev2026KeyRecovery}. Each returned vector has $N$
polynomial coefficients. Weis's construction instead solves for relations
among derivative conditions.

\begin{table}[htbp]
\centering\small\setlength{\tabcolsep}{4pt}
\caption{Primary parameters. All sets include infinity at multiplicity
six. The \(P_5\) lowered points are finite, retained, and non-unit.}
\label{tab:wopt-parameters}
\begin{tabular}{@{}lrrrrrrr@{}}
\toprule Parameter set & \(k\) & \(P\) & \(P_5\) & \(M\) & \(b\) & \(g\) & \(b'\)\\
\midrule
\texttt{mceliece348864} & 150 & 768 & 48 & 169 & 254 & 85 & 1016 \\
\texttt{mceliece460896} & 195 & 1248 & 3 & 213 & 320 & 107 & 1280 \\
\texttt{mceliece6688128} & 221 & 1664 & 41 & 239 & 359 & 120 & 1436 \\
\texttt{mceliece6960119} & 214 & 1547 & 30 & 232 & 348 & 116 & 1392 \\
\texttt{mceliece8192128} & 221 & 1664 & 41 & 239 & 359 & 120 & 1436 \\
\bottomrule
\end{tabular}
\end{table}

For each shortening, we choose the largest feasible $P_5$ subject to
$6(k+P)-P_5\geq7D_{\mathrm{sh}}+1-t$. Equality is feasible in every row:
\begin{equation}
 6(k+P)-P_5=7D_{\mathrm{sh}}+1-t,
 \qquad P_5=(15-m)t-k+6.
 \label{eq:wopt-reserve}
\end{equation}
The count includes the unit columns and the point at infinity. The
vanishing identity and the geometric rank conditions required for key
recovery are stated in \cite[Sec.~9]{Vedenev2026KeyRecovery}.
The sampler theorem and the resource bounds here hold for every matrix
with the specified derivative structure, irrespective of those conditions.

To choose $k$ and $P_5$, we use the mixed Hilbert-series prediction
\begin{equation}
 \kappa_{\mathrm{pred}}=[z^7]\bigl[(1+z)^k
             (1-z^2)^{P_6}(1-z^3)^{P_5}\bigr]_+,
 \label{eq:wopt-hilbert}
\end{equation}
where the series is truncated at its first nonpositive coefficient.
For each parameter set, we choose the smallest $k$ with a positive
prediction under \eqref{eq:wopt-reserve}. It gives, in parameter order,
\[
 965{,}303{,}280,\quad 9{,}469{,}314{,}720,\quad
 4{,}958{,}414{,}624,\quad 6{,}446{,}238{,}332,\quad 4{,}958{,}414{,}624.
\]
All five predictions are positive.

\subsection{A fixed attempt}
\label{sec:wopt-fixed-attempt}
Larger projection blocks shorten the sequences used by the sampler,
but require more storage. We choose the two widths together and count
both their saving in matrix products and their storage cost.
Write $n=N'=N+64$, $J=J_{\calK}$, $T=L'$ and $q=L+T$ for the sampler
of \Cref{subsec:wiedemann-sampler}. We use binary Toeplitz compression,
with $N+J+63$ independent random diagonal bits, and choose
\[
 b=\ceil{3M/2},\qquad g=b-M,\qquad b'=4b,\qquad
 L=\ceil{(n+64)/b},\qquad T=\ceil{(n+64)/b'}.
\]
Here $g\geq85$. Substitution in \Cref{thm:wiedemann-sampler} gives
\[
 \varepsilon
 \leq3\,2^{-64}+2^{-g}+\frac1{2^{b'-b-1}-1}<2^{-62}.
\]
Consequently independent retries multiply the expected cost by less than
$1+2^{-61}$. This correction is too small to change the displayed exponents.

The relation is computed by the circuit in \Cref{sec:wopt-generator}.
For reconstruction, retaining $U,W$ is enough; the initial block $[AU,W]$
need not remain stored. Split each relation coefficient into its first
$M$ rows and its last $g$ rows,
$C_j=[C_j^{\mathrm{top}};C_j^{\mathrm{bottom}}]$. Horner's rule then takes
the form
\begin{equation}\label{eq:wopt-reconstruction}
 H_L=UC_L^{\mathrm{top}},\qquad
 H_i=AH_{i+1}+UC_i^{\mathrm{top}}+WC_{i+1}^{\mathrm{bottom}}
 \quad\text{for }i=L-1,\ldots,0.
\end{equation}
Expanding $H_0$ gives $U+\sum_{j=1}^L A^{j-1}[AU,W]C_j$, which is the
output $U+Z$ of \Cref{subsec:wiedemann-sampler}. We test $K_+H_0=0$ before
discarding the padding coordinates.

Every attempt executes the same circuit. If relation extraction fails,
set all its coefficients to zero and retain a failure flag. Reconstruction
and every residual test still run. The flag prevents acceptance.

This schedule uses
\[
 Q_A=b(q-1)+M(L+1)\quad\text{applications of }A,
 \qquad M\quad\text{additional applications of }K_+.
\]
This count includes the $M$ products forming $AU$ and $ML$ products in
reconstruction. It is about $1.92n$ at all five parameter sets. The
reconstruction above uses $M$ more applications of $A$ than the Horner
formula in \Cref{subsec:wiedemann-sampler}; those applications are
included in $Q_A$.

\subsection{Where the gates are spent}
\label{sec:wopt-cost-decomposition}
Each application of $A$ first applies the hold-out matrix $K_+$ and then
the Toeplitz compression. Let $C_K$ and $C_T$ be the gate costs of these
two products. The Toeplitz seed is transformed once for sequence generation
and again for reconstruction, at cost $C_{\mathrm{pre}}$ each time.
The final check adds $M$ applications of $K_+$.

The other large operations are the projections $XA^iV$, reconstruction,
and relation computation. A dense binary product is charged for every
AND and XOR. This gives $2nb'bq$ gates for the projections and
$2nbML+nM$ for reconstruction, including the final addition of $U$.
Let $G$ include constructing the approximant basis and extracting the
relation. There are $S_F=bM(L+1)$ relation bits to mask if extraction
fails. With the residual tests and acceptance flag, the total is
\begin{equation}
 \begin{aligned}
 C_1={}&Q_A(C_K+C_T)+MC_K+2C_{\mathrm{pre}}\\
       &+2nb' bq+2nbML+nM+G+S_F+3JM+4 .
 \end{aligned}
 \label{eq:wopt-totalgates}
\end{equation}

The first line counts the repeated matrix products and seed preparations.
The second counts projections, reconstruction, the relation, and the
final checks. For \texttt{mceliece348864}, the products with $K_+$ account
for $33.7\%$ of the total and the Toeplitz products and preparations for
$43.7\%$. The other operations account for $22.6\%$.

\subsection{Sparse products and the mixed derivative map}
\label{sec:wopt-sparse}
We first evaluate orders five and three at every retained point, using
the grouped product of \Cref{sec:wopt-groupcount}. Call this common
operator $K_6$, with $J_6=P(\Cfive+\Cthree)$ rows. At a point lowered to
multiplicity five, we replace the order-five values by order-four values.
No new polynomial evaluation is needed: for a homogeneous multilinear
polynomial of degree seven,
\begin{equation}
 D_T^{[4]}f(y)=\sum_{u\notin T}y_uD_{T\cup\{u\}}^{[5]}f(y).
 \label{eq:wopt-euler}
\end{equation}
Each degree-three monomial on the right occurs three times, hence once
over $\FF_2$. The order-three derivative values themselves are retained.
This gives a public linear map from $K_6v$ to $Kv$.

The priced schedule first forms the full \(J_6\)-bit output and then
applies \eqref{eq:wopt-euler}. When \(P_5>0\), it reserves a separate
\(J_6\)-bit temporary. It charges the map
\[
 C_{\mathrm{mix}}=10P_5\Cfive
\]
gates and at most
\[
 V_{\mathrm{mix}}=
 \begin{cases}
 J_6+2J+10P_5\Cfive,&P_5>0,\\
 0,&P_5=0
 \end{cases}
\]
bit transfers. The gate term is zero when \(P_5=0\).
The $K_6$ output remains stored until this conversion is complete;
its temporary storage and transfers are included in the bounds.

\subsubsection{Sharing work between derivative rows}
\label{sec:wopt-groupcount}
The order-five rows evaluate many quadratic forms at the same public
points. We collect evaluations that share a right operand into matrix
products. This is where fast matrix multiplication reduces the gate count.

A scalar input is the coefficient vector $(h_I)_{|I|=7}$.
For a five-set $S$, the order-five derivative at the public point $y_j$ is
\[
 \sum_{\substack{u<v\\\{u,v\}\cap S=\varnothing}}
 h_{S\cup\{u,v\}}y_{j,u}y_{j,v}.
\]
The order-three rows need partial sums from these evaluations. For a
pair $V\subset S$, the required sum uses $u>\max V$. Write
$S=\{s_0<s_1<s_2<s_3<s_4\}$. The possible values of $\max V$ are
$s_1,s_2,s_3,s_4$, so these are exactly the cuts we need.
Partition the pair coordinates by the five intervals between
$-1,s_1,s_2,s_3,s_4,k$. For each interval $(a,c)$, collect all five-sets
having that interval.
The matrix with entries $h_{S\cup\{u,v\}}$ for $a<u\leq c$, $u<v$,
multiplied by the matrix with entries $y_{j,u}y_{j,v}$, gives all their
segment sums. An entry with $S\cap\{u,v\}\ne\varnothing$ is a constant
zero. Thus many derivative rows share the same public right operand.

Each segment contributes to the order-five row. It also contributes
$y_{j,V}$ times its value to the order-three row $S\setminus V$ for each
pair $V\subset S$ with $\max V\leq a$. To verify this, fix a three-set
and four further indices. There is exactly one split of those four into
the two smaller indices $V$ and two larger indices $u,v$.
The segment condition selects precisely this split. The five intervals
have $0,1,3,6,10$ eligible pairs, giving twenty lower-order emissions
per five-set. No order-three condition is omitted.

The size of each matrix product follows by counting its five-sets and
pairs. For an interior interval, both endpoints belong to $S$; $j$
elements lie below $a$ and $3-j$ lie above $c$, with $1\leq j\leq3$.
The numbers of rows and pair columns are
\[
 \begin{aligned}
 R_{-1,c}&=c\binom{k-c-1}{3},&R_{a,k}&=\binom a4,\\
 R_{a,c}&=\sum_{j=1}^{3}\binom aj\binom{k-c-1}{3-j},&
 H_{a,c}&=\binom{k-a-1}{2}-\binom{\max(k-c-1,0)}2.
 \end{aligned}
\]
Invalid binomial coefficients are zero. Counting the five intervals of
each five-set gives $\sum R_{a,c}=5\Cfive$, and partitioning all its pair
columns gives $\sum R_{a,c}H_{a,c}=\Cfive\binom{k}{2}$.

\subsubsection{Counting the grouped products}
Each group is divided into blocks of at most $s=2048$ rows and pair
coordinates, retaining all $P$ point columns. We use the following
seven-product Winograd step within these blocks. Its four input additions
on each side and seven output additions determine the constants below.
For block matrices $A,B$, all signs are plus over $\FF_2$:

\[
\begin{array}{ll}
 S_1=A_{21}+A_{22},&T_1=B_{12}+B_{11},\\
 S_2=S_1+A_{11},&T_2=B_{22}+T_1,\\
 S_3=A_{11}+A_{21},&T_3=B_{22}+B_{12},\\
 S_4=A_{12}+S_2,&T_4=T_2+B_{21}.
\end{array}
\]
Its seven recursive products are
\[
\begin{aligned}
 P_1&=A_{11}B_{11},&P_2&=A_{12}B_{21},&P_3&=S_4B_{22},\\
 P_4&=A_{22}T_4,&P_5&=S_1T_1,&P_6&=S_2T_2,&P_7&=S_3T_3.
\end{aligned}
\]
Put $U_2=P_1+P_6$, $U_3=U_2+P_7$ and $U_4=U_2+P_5$.
The four output blocks are $P_1+P_2$, $U_4+P_3$, $U_3+P_4$ and
$U_3+P_5$, which expand to $AB$.
There are four left-operand additions, four public-operand additions and
seven output additions. Reversing the circuit linear in the left operand
exchanges the four and seven counts. These recurrences give the following
closed costs.

A depth-\(e\) Winograd product on dimensions \(a,h,c\),
padded to multiples \(a',h',c'\) of \(u=2^e\), costs
\[
 \begin{split}
 F_e&=7^e(a'/u)(c'/u)(2h'/u-1)
       +\frac{(4a'h'+7a'c')(7^e-4^e)}{3\,4^e},\\
 T_e&=F_e+a'(c'-h'),\qquad
 U_e=\frac{4h'c'(7^e-4^e)}{3\,4^e}.
 \end{split}
\]
Here \(F_e\) is the variable forward cost, \(T_e\) the corresponding
transposed cost,
and \(U_e\) the public preparation cost. For each block, choose
\(0\le e\le11\) minimizing \((F_e+T_e,U_e,e)\) lexicographically.
Only \(F_e\) is used in this sampler. For each block of pair coordinates, a public
preparation is shared by all row blocks using that depth and is otherwise
rebuilt. At the ends of the recursion there are $7^e$ public right operands,
each of size $(h'/u)\times(c'/u)$. Retaining them takes
$7^e(h'/u)(c'/u)$ bits.

Let \(F_{\mathrm{grp}}\) be the sum of all forward block costs, including
\(R_{a,c}P\ceil{H_{a,c}/s}\) gates for partial-product accumulation in
each group. Let \(U_{\mathrm{grp}}\) sum the distinct public preparations
and the \(PH_{a,c}\) pair-mask gates of each group. Then
\begin{equation}
 C_K^{\mathrm{grp}}=F_{\mathrm{grp}}+U_{\mathrm{grp}}+65P\Cfive+C_{\mathrm{mix}},
 \qquad
 V_K^{\mathrm{grp}}=21N+50P\Cfive+J_6+V_{\mathrm{mix}}.
 \label{eq:wopt-groupedK}
\end{equation}
The \(65P\Cfive\) gates consist of five top-row accumulations and
twenty lower-order emissions per point and five-set. Each emission uses
one public-mask AND, one AND with the segment value and one XOR.
The transfer bound counts \(21N\) coefficient incidences, read/write pairs
for the five top and twenty lower accumulations, and output initialization. Public preparations are charged
again on every scalar application of \(K\).

Let \(L_{\mathrm{leaf}}\) be the maximum sum of retained public leaf bits
over one group, and let \(\bar P=s\ceil{P/s}\). A conservative local
region is
\begin{equation}
 \begin{split}
 L_K^{\mathrm{grp}}={}&L_{\mathrm{leaf}}
 +s\min\{s,\tbinom{k}{2}\}+sP\\
 &+16s(s+\bar P)+8s\bar P+Pk+2^{20}.
 \end{split}
 \label{eq:wopt-groupcache}
\end{equation}
All scratch is reused between groups. The two recursive scratch terms
allow depth-first evaluation of the seven products and their preparations;
at each level the operand dimensions halve, so the live matrix areas form
a geometric sum. The factors $16$ and $8$ exceed that sum for the variable
and public arrays, respectively. The retained leaf forms are counted
separately, including every depth used by a row block.

\subsection{Toeplitz multiplication and static pruning}
\label{sec:wopt-fft}
A Toeplitz product is a segment of the coefficient vector of a binary
polynomial product. We compute it by evaluation and interpolation:
transform the diagonal seed and the input, multiply their values at the
same points, and interpolate. The seed transform is reused across all
products in a phase. Each new input therefore needs one forward transform,
pointwise multiplication, and one inverse transform.

Let $S_R$ be the seed length, $c_s,c_x$ the numbers of packed coefficients,
and $Q$ the transform length:
\[
 S_R=n+J-1,\qquad c_s=\ceil{S_R/64},\qquad c_x=\ceil{J/64},
 \qquad Q=2^\kappa\ge c_s+c_x,
\]
with $Q$ the least such power of two. To explain the packing, let
$\mathbb E=\FF_2[z]/(z^{128}+z^7+z^2+z+1)$. Pack $64$ binary coefficients
into one polynomial in $z$ of degree below $64$. Products of two chunks
have degree below $127$, so multiplication in $\mathbb E$ gives their
ordinary binary product. After interpolation, overlapping chunks are
added with offsets of $64$ bits. The desired Toeplitz output is the middle
$n$ bits of this binary product.

Cantor's additive transform evaluates on a binary vector space of field
elements \cite{Cantor1989}. Its basis is chosen so that polynomial
reduction splits an evaluation into two smaller evaluations. Take $\beta_0=1$ and define $\beta_i$ by
$\beta_i^2+\beta_i=\beta_{i-1}$, choosing the solution with zero constant
coefficient in the polynomial basis of $\mathbb E$. Let
$\beta(h)=\sum_j h_j\beta_j$ for the binary digits $h_j$ of $h$.
For all depths used here ($\kappa\leq40$), the required basis elements
lie in the $64$-dimensional subfield.
The associated subspace polynomials satisfy $\mathfrak s_0(X)=X$ and
$\mathfrak s_{i+1}(X)=\mathfrak s_i(X)^2+\mathfrak s_i(X)$; explicitly,
$\mathfrak s_i(X)=\sum_{j\subseteq i}X^{2^j}$, where $j\subseteq i$ means that
every set bit of $j$ is a set bit of $i$.

The subspace polynomials give the following recursive evaluation.
Each node divides a polynomial into two remainders, one for each coset. At level $i$ a node
has $2h$ coefficients, $h=2^i$, and fixed multiplier (twiddle) $\beta(2v)$, where $v$ is its
index among the nodes at that level. Visit $t=2h-1,\ldots,h$ in that
order. For every proper bit subset $j\subsetneq i$, add coefficient $t$
to coefficient $t-h+2^j$. After these additions, add $\beta(2v)$ times
coefficient $t$ once to coefficient $t-h$. Then add each lower-half coefficient to the matching
upper-half coefficient. Recursing on both halves evaluates the polynomial
on the two cosets, since their subspace-polynomial values differ by one.
The inverse undoes these updates in reverse order. In particular, the
monomial divisions are included in the transform cost.

\subsubsection{The field operations used by the FFT}
The multiplications inside the transforms use fixed elements of the
$64$-dimensional subfield. Representing a field
element by two coordinates over that subfield lets such a multiplication
act separately on the coordinates. Only the pointwise products need a
general $128$-bit field multiplication.

We use the quadratic tower
\[
 \mathbb B=\FF_2[\eta]/(\eta^{64}+\eta^{63}+\eta^8+\eta+1),\qquad
 \mathbb E=\mathbb B[\theta]/(\theta^2+\theta+\eta).
\]
The isomorphism to the preceding polynomial representation sends $\eta$
and $\theta$ to the following bit strings, with the low bit the constant
coefficient:
\[
\begin{aligned}
 \eta&\longmapsto\texttt{865b53d9ffb9034f9099a9c58787d41a},\\
 \theta&\longmapsto\texttt{75c1f22587f712dc4a5cdc3b43e79fcc}.
\end{aligned}
\]
These specify both binary conversion matrices: the columns of the map
from the tower are $1,\eta,\ldots,\eta^{63},\theta,\theta\eta,\ldots,
\theta\eta^{63}$ in $\mathbb E$, and the reverse map is its binary inverse.
The defining identities can be checked by polynomial reduction. The two
base polynomials pass the irreducibility test
$X^{2^u}=X\pmod f$ and $\gcd(X^{2^{u/2}}-X,f)=1$ for $u=64,128$;
also $\operatorname{Tr}_{\mathbb B/\FF_2}(\eta)=1$, so the quadratic
extension is a field.

For a $v$-bit binary polynomial product, use schoolbook multiplication
at $v=4$ and the three-product recursion at larger sizes. Its gate count is
$R(4)=25$, $R(2v)=3R(v)+7v-3$, whence $R(64)=3725$.
Reduction modulo the pentanomial adds $252$ XORs.
For $a=a_0+\theta a_1$ and $c=c_0+\theta c_1$, compute
$z_0=a_0c_0$, $z_2=a_1c_1$ and $z_1=(a_0+a_1)(c_0+c_1)$.
Their product is $(z_0+\eta z_2)+\theta(z_1+z_0)$.
The additions use $256$ XORs and multiplication by $\eta$ another four.
Thus a variable field product costs $3(3725+252)+260=12191$ gates.

A binary linear map with row weights $w_j$ costs
$\sum_j\max(0,w_j-1)$ XORs. Applying this formula to the specified
conversion matrices gives $3919$ XORs to embed a $64$-bit chunk and
$8097$ to convert a full output coefficient. These are fixed maps;
the Toeplitz seed and both pointwise operands remain variable inputs.

For the fixed twiddle multiplication matrices, let \(a_d\) be the number
of active entries in the union of the first \(d\) nonconstant basis
matrices. Their verified values begin
\[
 a_d=0,2339,3381,3580,3788,3916,3984,4053,4096,\ldots .
\]
All $a_d$ with $d\geq8$ equal $4096$. Each multiplication matrix is
obtained by reducing $\beta_j\eta^t$ modulo the defining polynomial of
$\mathbb B$, for $0\leq t<64$.
Let $\operatorname{popcount}(i)$ be the number of ones in the binary
expansion of $i$, and put $s_i=2^{\operatorname{popcount}(i)}$.
Before removing any operations, a transform costs
\begin{equation}
 C_F=64Q\sum_{i=0}^{\kappa-1}s_i
          +\frac Q2\sum_{d=1}^{\kappa-1}a_d.
 \label{eq:wopt-fulltransform}
\end{equation}
At level $i$ there are $2^{\operatorname{popcount}(i)}-1$ lower
subspace-polynomial terms and one half-addition. The twiddle matrices
range over a binary span: each nonzero matrix entry is one in exactly
half its elements. Summing these two contributions gives the formula.

\subsubsection{Removing operations on known zeros and unused outputs}
The packed input does not fill the transform array, and the Toeplitz
product uses only a middle segment of the interpolated product. We remove
operations whose inputs are known to be zero from this padding, and
operations that cannot affect the required output. These decisions depend
only on the dimensions; they do not inspect random data.

For a forward input supported on its first $c$ field coefficients,
let $F_\kappa(c)$ be the resulting gate count. If \(c\le1\), \(F_\kappa(c)=0\).
Otherwise put \(v=\ceil{\log_2 c}\), \(u=c-2^{v-1}\) and
\(e=\kappa-v\). Then
\begin{equation}
 \begin{split}
 F_\kappa(c)={}&64Q\sum_{i=0}^{v-2}s_i
             +128\,2^e u\,s_{v-1}\\
             &+\frac Q2\sum_{d=e+1}^{\kappa-1}a_d
             +2^e u\,a_e .
 \end{split}
 \label{eq:wopt-forwardprune}
\end{equation}
Empty sums are zero. A copied polynomial still has to be evaluated
on each of its cosets, so the work below every copy is retained.

The inverse must supply the inclusive interval \([a,z]\), where
\[
 a=\max\{0,\lfloor(J-1)/64\rfloor-1\},\qquad
 z=\lfloor(J+n-2)/64\rfloor,\qquad c_o=z-a+1.
\]
The interval's position matters, not just its length. To count the
inverse gates needed for it, mark all $128$ bits of every output
coefficient in $[a,z]$ and work backward through the inverse graph.
Call the resulting count $I_\kappa[a,z]$.
For each update \(d\leftarrow d+\beta s\), visited in reverse execution
order, retain the update exactly when \(d\) is marked, and then mark
\(s\). A retained field XOR costs 128 gates; a retained nonzero twiddle
update costs twice the weight of its 64-bit multiplication matrix.
The inverse graph is the forward graph in reverse order, so this rule
fully specifies a finite parameter-only circuit.

The forward count follows the same recursion after removing structurally
zero inputs. When the input occupies more than half a block, its $u$
upper coefficients activate the corresponding updates in each of the
$2^e$ copied subproblems. This gives the displayed forward formula.
For the inverse, the backward marking rule itself defines the retained
finite graph and its count. No test of the values of
random or intermediate coefficients removes a gate.

The Toeplitz costs are therefore
\begin{equation}
 \begin{split}
 C_{\mathrm{pre}}&=F_\kappa(c_s)+3919c_s,\\
 C_T&=F_\kappa(c_x)+I_\kappa[a,z]+12191Q
           +3919c_x+8097c_o+n .
 \end{split}
 \label{eq:wopt-toeplitzcost}
\end{equation}
There are two seed preparations per attempt. All word-transfer
allowances below remain those of the unpruned graph: wire copies are
not made free by deleting their XOR gates.

\subsection{Computing the relation with fixed addresses}
\label{sec:wopt-generator}
\Cref{subsec:wiedemann-relation} obtains the required relation from a
reduced approximant basis.
For a circuit count, choosing a pivot or reading a coefficient at a
computed degree must itself be implemented by gates. We keep the large
coefficient arrays in fixed positions and apply all pivot choices through
a small matrix of control bits.

Put \(p=b+b'\), \(d_q=\ceil{\log_2(q+2)}\),
\(d_p=\ceil{\log_2p}\), \(w=d_q+d_p\), and
\(H_{\triangle}=q(q+1)/2\). We build the basis one order at a time, with the shift from
\Cref{subsec:wiedemann-relation}. Store its top coefficient matrices
$D_j\in\FF_2^{b\times p}$, $0\leq j\leq q$. Of the bottom polynomial
matrix, only its current coefficient $P_i\in\FF_2^{b'\times p}$ is needed
for the next equation.
The \emph{discrepancy} is the coefficient of order $i$ that must be
cancelled to impose the next approximation equation. It is
\[
 \Delta_i=\sum_{j=0}^i S_{i-j}D_j+P_i.
\]
Start with the identity basis. Order its columns by shifted degree and
leading position. Eliminate the discrepancy in that order, adding earlier
pivot columns to later columns, and multiply every pivot column by the
polynomial variable. This imposes the next coefficient equation. Choosing
pivots in shifted-degree order preserves independent shifted leading
vectors, so induction gives a reduced approximant basis. The extraction
lemma, \Cref{lem:wiedemann-predictable}, then applies.

\paragraph{Implementing pivot choices.}
Sort the small discrepancy and control arrays with an ascending odd-even merge network
\cite{Batcher1968}. Pad the network to a power of two with infinite keys
and omit comparisons whose larger index is outside the $p$ real columns.
An infinite key never moves to a lower index. Record every conditional
swap and elimination as a column operation on a small matrix $E_i$,
initially $I_p$. To find a pivot row, scan all candidate rows and set
only the flag of the first eligible row. Every possible update is then
performed with its flag as an AND mask. This represents the chosen column
operations in $E_i$ while keeping all addresses fixed.

Let $D_i^{\mathrm{piv}}$ be the diagonal matrix of pivot flags. For each
coefficient index, in descending order, compute
\[
 B_j=D_jE_i,\qquad
 D'_j=B_j(I-D_i^{\mathrm{piv}})+B_{j-1}D_i^{\mathrm{piv}},\qquad B_{-1}=0.
\]
Three $b\times p$ coefficient buffers suffice for this descending pass.
The next bottom tail is $(P_iE_i)D_i^{\mathrm{piv}}$. Before the update the
bottom polynomials have degree at most $i$, so their discarded
coefficients can never enter a later discrepancy.

\paragraph{Extracting the prescribed constant coefficient.}
The extraction lemma needs the coefficient at a column's shifted degree.
Scan all coefficient positions up to $L$ and mask the
coefficient whose index equals a column's shifted degree. Solve the
resulting $b\times p$ binary systems for the prescribed columns of $E$.
Full pivot scans, controlled swaps and controlled elimination cost at most
$64p^4$ gates, including the inconsistency test. Finally reverse one top
basis column into an array of length $q+L+1$. For a column of shifted
degree $\rho_j\leq L$, shift it toward the beginning by $q-\rho_j$
using fixed stages of offsets $1,2,4,\ldots$. Each stage selects
between two fixed addresses. Mask columns of degree greater than $L$ and
accumulate the selected columns into all $M$ relations. This implements
\Cref{lem:wiedemann-predictable} without a degree-indexed read.

\paragraph{Gate count.}
Let $s_p$ be the number of comparisons left after the dummy comparisons
are removed. \Cref{tab:wopt-generator-gates} separates the work into
basis construction and extraction. Their sums give
\begin{equation}\label{eq:wopt-generator}
 G=G_{\mathrm{basis}}+G_{\mathrm{extract}}.
\end{equation}
A conditional bit swap costs four gates, a bit selection three, and a
comparison at most eight per key bit. The table charges every operation,
including those whose control bit is zero and those on a failed attempt.

\begin{table}[htbp]\centering\small
\caption{Gates in the relation circuit. The first seven rows sum to
$G_{\mathrm{basis}}$; the last five sum to $G_{\mathrm{extract}}$.
$H_{\triangle}=q(q+1)/2$ counts the coefficient positions visited over all
orders.}
\label{tab:wopt-generator-gates}
\begin{tabular}{@{}ll@{}}
\toprule Operation & Gates\\\midrule
Discrepancies & $2b'bpH_{\triangle}$\\
Sorting & $qs_p(12w+4b'+4p)$\\
Pivot selection and elimination & $q[p(3b'+1)+\tfrac{p(p-1)}2(4b'+2p-1)]$\\
Top-coefficient transformations & $2bp^2H_{\triangle}$\\
Selected coefficient shifts & $3bp(H_{\triangle}+q)$\\
Bottom-tail updates & $q(2b'p^2+b'p)$\\
Degree increments & $2d_qpq$\\\midrule
Leading-coefficient scans & $p(L+1)(2d_q+2b)$\\
Binary solve and consistency test & $64p^4$\\
Extraction control & $p(24d_q+M+4)$\\
Fixed-stage shifts & $3pb(q+L+1)d_q$\\
Accumulating the relations & $2pbM(L+1)$\\\bottomrule
\end{tabular}
\end{table}

\subsection{What must remain in memory}
\label{sec:wopt-storage}
Storage peaks while the approximant basis is being formed. The full
projected sequence and the top basis rectangle are both live then.
Together they occupy about $92\%$ of the peak in the first parameter set.
The FFT arrays have already been released, so their capacity
is not added to this phase.

All capacities below are in bits. The projected sequence, top basis
rectangle and temporary array for relation extraction occupy
\[
 S_{\mathrm{seq}}=b' bq,\qquad S_D=bp(q+1),\qquad S_H=b(q+L+1).
\]
The relation itself uses the $S_F$ bits defined in
\Cref{sec:wopt-cost-decomposition}. Field arithmetic uses a workspace of
$2^{20}$ bits. A further $2^{12}$ field coefficients are kept in a cache
while the smallest FFT subproblems are evaluated. Their combined capacity
is $W_c=3\cdot2^{19}$ bits.

The local workspaces for projection, reconstruction and relation control
are $L_X,L_H,L_G$. The common stored data use $B_0$ bits, and $B_{\mathrm{mix}}$
is the extra derivative buffer at lowered points. The operator bank,
including its two transform arrays, has capacity $W_A$:
\[
 \begin{aligned}
 L_X&=b' b+b'+b+4096,&
 L_H&=bM+b+2M+4096,\\
 L_G&=16p^2+4pw+b' b+3bp+2s_pd_p+4096,\\
 B_0&=S_R+J+Pk+W_c,\qquad&
 B_{\mathrm{mix}}&=J_6\,\mathbf1_{\{P_5>0\}},\\
 W_A&=256Q+2n+B_0+L_K-Pk+B_{\mathrm{mix}}.
 \end{aligned}
\]
Here $L_K=L_K^{\mathrm{grp}}$ is the region in
\eqref{eq:wopt-groupcache}. The subtraction of $Pk$ avoids counting the
public-point copy twice. The term $256Q$ holds two arrays of $Q$ field
elements, and $2n$ allows the scalar input and output buffers.

Sequence generation keeps $U,W$, the current block, and $X$. Basis
construction and extraction retain only $U,W$ of these long blocks.
Reconstruction adds the current $M$ output columns. Including the
coefficient arrays and local workspaces gives
\begin{equation}
 \begin{split}
 W_{\mathrm{seq}}&=n(2b+b')+S_{\mathrm{seq}}+W_A+L_X,\\
 W_{\mathrm{gen}}&=nb+S_{\mathrm{seq}}+S_D+B_0+L_G,\\
 W_{\mathrm{ext}}&=nb+S_D+S_F+S_H+B_0+L_G,\\
 W_{\mathrm{rec}}&=n(b+M)+S_F+W_A+L_H .
 \end{split}
 \label{eq:wopt-phase}
\end{equation}
The full mixed-output temporary and the large transform arrays are released
before the generator. The seed remains stored, and its transform is
computed again for reconstruction; both preparations are counted. The generator determines the peak in every
displayed one-attempt configuration.

\subsection{FFT locality with all workspace movement charged}
\label{sec:wopt-locality}
The gate count alone does not charge for bringing a coefficient to the
arithmetic workspace. Following \cite[Section 6]{Weis2026}, we charge
operations in a region of $x$ bits by $\rho(x)=\ceil{\sqrt{x}}$.
A gate uses its local region. A transferred bit uses a region containing
both its source and destination. Thus the total is
\begin{equation}
 \mathcal A_1=\sum_u C_u\,\rhoo{L_u}
                  +\sum_v V_v\,\rhoo{W_v},
 \qquad \rho(x)=\ceil{\sqrt{x}}.
 \label{eq:wopt-accessmodel}
\end{equation}
To use a small region for an FFT node, the arithmetic workspace must be
beside that node's coefficients. We move the workspace as the transform
proceeds and charge these movements as well as the operand transfers.

\subsubsection{One movable workspace}
Reserve a block of $W_c$ bits between the two transform arrays for the
workspace. This is the ``gap'' in the layout below; it holds the same
workspace throughout the computation. For a transform level \(i\ge r=12\), the active
coefficient blocks have size
\[
 B_i=2^{i+8},\qquad n_i=Q/2^{i+1}.
\]
Evaluate the nodes of one level from left to right, and those of the
next level from right to left. Nodes at the same level are independent,
so this change of order does not change the transform. After computing a node beside the gap,
rotate \([\text{gap}][\text{node}]\) into
\([\text{node}][\text{gap}]\), or the reverse. The next node is adjacent
to the same workspace. A fixed cycle permutation uses bounded scratch
and at most \(2(B_i+W_c)\) bit transfers within that union. No cache is
replicated per node.

The level's original operand traffic is
\begin{equation}
 V_i=(128s_i+256)Q-256\,2^i.
 \label{eq:wopt-leveltraffic}
\end{equation}
At each level exactly one node has zero twiddle and needs no twiddle
operand transfers. This accounts for the subtracted term. Let
\(B_r^{\mathrm{res}}=128\,2^r\). A transform's nonresident-node charge,
including workspace movement, is bounded by
\begin{equation}
 \begin{split}
 \Xi={}&\sum_{i=r}^{\kappa-1}
       [V_i+2n_i(B_i+W_c)]\,\rhoo{B_i+W_c}\\
 &+2(Q/2^r)(B_r^{\mathrm{res}}+W_c)\,
                      \rhoo{B_r^{\mathrm{res}}+W_c}\\
 &+\eta\,2(128Q+W_c)\rhoo{W_A},
 \qquad
 \eta=\mathbf1_{\{\kappa-r+1\ {\mathrm{odd}}\}} .
 \end{split}
 \label{eq:wopt-xi}
\end{equation}
The second line moves the gap past the resident subtrees. The final
line returns it to the starting endpoint when the number of scans is
odd. Both forward and inverse transforms then preserve the surrounding
bank layout. The existing arrays and single workspace suffice.

A subproblem of $2^r$ coefficients fits in the cache. Once it has been
loaded, its bottom $r$ levels run there. Summing their operand transfers
over all cosets gives
\[
 V_{\mathrm{low}}
 =128Q\sum_{i=0}^{r-1}s_i+256rQ-256(2^r-1).
\]
It is charged in \(W_c\), the union of resident data and field workspace.
Loading and storing all resident subtrees costs another \(256Q\) bank
bits per transform; this charge is retained.

\subsubsection{Combining arithmetic and transfers}
Packing the input, loading and storing cached subproblems, multiplying
values pointwise, and reading out the middle window use the full operator
bank. Their transfer counts are
\[
 \begin{split}
 V_{\mathrm{pre}}^{\mathrm{outer}}&=S_R+128Q+256Q,\\
 V_T^{\mathrm{outer}}&=J+128Q+512Q+384Q+128c_o+n.
 \end{split}
\]
The $128Q$ terms include storing padded zero coefficients. The pointwise
products read both operands and write the result. Combining these passes
with local arithmetic, cached transfers and workspace movement gives
\begin{equation}
 \begin{split}
 \mathcal T_{\mathrm{pre}}
 &=C_{\mathrm{pre}}\rhoo{2^{20}}+V_{\mathrm{low}}\rhoo{W_c}
       +V_{\mathrm{pre}}^{\mathrm{outer}}\rhoo{W_A}+\Xi,\\
 \mathcal T
 &=C_T\rhoo{2^{20}}+2V_{\mathrm{low}}\rhoo{W_c}
       +V_T^{\mathrm{outer}}\rhoo{W_A}+2\Xi.
 \end{split}
 \label{eq:wopt-toeplitzaccess}
\end{equation}
There are \(Q_A\) actions and two preparations. The bank capacity
contains the mixed temporary and the selected sparse workspace even
when those arrays are inactive during a transform.

The other parts of the sampler are charged in the same way. Each of the $Q_A+M$ applications of $K$
costs $C_K\rho(L_K)+V_K\rho(W_A)$. Add the local gate charges
\[
 \begin{split}
 &2nb' bq\,\rho(L_X)+(2nbML+nM)\rho(L_H)\\
 &\qquad +(G+S_F+1)\rho(L_G)+(3JM+3)\rho(4096).
 \end{split}
\]
Here $C_K,V_K,L_K$ are the grouped choices. The remaining transfers are
listed in \Cref{tab:wopt-other-transfers}; each count is multiplied by the
square root, rounded upward, of the indicated region size.
Define
\[
\begin{split}
 R_G={}&H_{\triangle}(b' b+3bp)+qbp+S_D+bp(L+1),\\
 R_E={}&p\bigl[2b(q+1)+S_H+3S_Hd_q+(b+2bM)(L+1)\bigr]+S_F.
\end{split}
\]
The terms of $R_G$ are sequence reads, coefficient reads and writes,
initialization, and leading-coefficient scans. Those of $R_E$ load and
clear the shift array, execute its stages, accumulate the output, and
initialize that output.

\begin{table}[htbp]\centering\small
\caption{Transfers in addition to the sparse and Toeplitz products.
$W_*=\max W_{\mathrm{phase}}$. Every row is charged even on extraction failure.}
\label{tab:wopt-other-transfers}
\begin{tabular}{@{}p{.27\linewidth}p{.49\linewidth}l@{}}
\toprule Operation & Bit transfers & Region\\\midrule
Input stores and helpers & $n(b+b')+S_R+Pk+2ng$ & $W_{\mathrm{seq}}$\\
Sequence column copies & $4n[M+b(q-1)]$ & $W_{\mathrm{seq}}$\\
Projection operands & $q[n(b+b')+b' b]$ & $W_{\mathrm{seq}}$\\
Projection resets & $qb' b$ & $L_X$\\
Basis and leading scans & $R_G$ & $W_{\mathrm{gen}}$\\
Controller stores & $q(p^2+2b' p)+4L_G$ & $L_G$\\
Extraction and guard & $R_E+2S_F$ & $W_{\mathrm{ext}}$\\
Recovery column copies & $4nML$ & $W_{\mathrm{rec}}$\\
Reconstruction & $(L+1)[n(b+2M)+bM]$ & $W_{\mathrm{rec}}$\\
Final input copies & $2nM$ & $W_{\mathrm{rec}}$\\
Residual reads & $JM$ & $W_A$\\
Returned tuple and flag & $NM+1$ & $W_{\mathrm{rec}}$\\
Three layout changes & $6W_*$ & $W_*$\\\bottomrule
\end{tabular}
\end{table}

The phase allocations use compact layouts. At each boundary a fixed
permutation of the largest arena moves the surviving arrays to their new
positions. Evaluating its cycles with one-bit scratch reads and writes
each position at most once, giving the last row of the table.
The gate charges, the sparse charges, $Q_A\mathcal T+2\mathcal T_{\mathrm{pre}}$,
and this transfer table together give the access column of
\Cref{tab:wopt-optimized-comparison}.

\printbibliography[heading=bibintoc]
\end{document}
